\documentclass[11pt,a4paper]{article}
\usepackage[utf8]{inputenc}
\usepackage[T1]{fontenc}
\usepackage[spanish,es-minimal,es-noshorthands]{babel}
\AtBeginDocument{\renewcommand{\sin}{\operatorname{sen}}\renewcommand{\arcsin}{\operatorname{arcsen}}}
\usepackage{lmodern}
\usepackage[sfdefault,lining]{FiraSans}
\usepackage[a4paper,top=2.3cm,bottom=2.2cm,left=1.8cm,right=1.8cm,headheight=1.3cm,headsep=0.4cm,footskip=1cm]{geometry}
\usepackage{amsmath,amssymb}
\usepackage{icomma}
\usepackage{graphicx}
\usepackage[table]{xcolor}
\usepackage{tikz}
\usepackage{fontawesome5}
\usepackage{ccicons}
\usepackage{fancyhdr}
\usepackage{lastpage}
\usepackage{tasks}
\usepackage{needspace}
\usepackage[most]{tcolorbox}
\usepackage{array}
\usepackage[hidelinks]{hyperref}
\AtBeginDocument{\hypersetup{pdftitle={Ficha de ejercicios de la unidad 1: Números naturales y enteros. 2º ESO},pdfsubject={Matemáticas 2º ESO},pdfkeywords={números naturales, división entera, múltiplos, divisores, números primos, mcm, MCD, números enteros, operaciones combinadas, potencias, raíces}}}

\definecolor{navy}{HTML}{1B2748}
\definecolor{accent}{HTML}{E8603A}
\definecolor{teal}{HTML}{1F9E90}
\definecolor{soft}{HTML}{EEF1F8}
\definecolor{muted}{HTML}{5D6685}

\setlength{\parindent}{0pt}
\setlength{\parskip}{3pt}
\settasks{label=\alph*), label-format=\bfseries\color{navy}, label-width=1.3em, item-indent=1.8em, column-sep=0.6em, after-item-skip=2pt, before-skip=3pt, after-skip=4pt}

% ---------- capçalera i peu ----------
\newcommand{\llicencia}{\href{https://creativecommons.org/licenses/by-nc-sa/4.0/deed.es}{\ccbyncsa}\ \ Licencia Creative Commons Atribución-NoComercial-CompartirIgual 4.0 Internacional (CC BY-NC-SA 4.0)}
\pagestyle{fancy}
\fancyhf{}
\fancyhead[L]{\raisebox{-0.35\height}{\includegraphics[height=0.95cm]{logo-gva.png}}}
\fancyhead[C]{\small\color{muted}Ficha de ejercicios · Unidad 1: Números naturales y enteros · 2\textsuperscript{o} ESO}
\fancyhead[R]{\raisebox{-0.35\height}{\includegraphics[height=0.95cm]{logo-vermellar.png}}}
\fancyfoot[L]{\scriptsize\color{muted}\llicencia}
\fancyfoot[R]{\small\color{muted}\thepage}
\renewcommand{\headrulewidth}{0.4pt}
\renewcommand{\headrule}{\color{muted!40}\hrule height \headrulewidth}
\fancypagestyle{primera}{\fancyhf{}\renewcommand{\headrulewidth}{0pt}
  \fancyfoot[L]{\scriptsize\color{muted}\llicencia}
  \fancyfoot[R]{\small\color{muted}\thepage}}

% ---------- nivells i exercicis ----------
\newcommand{\estrelles}[1]{{\color{accent}\ifcase#1\or\faStar\faStar[regular]\faStar[regular]\or\faStar\faStar\faStar[regular]\or\faStar\faStar\faStar\fi}}
\newcounter{exer}
\newcommand{\ex}[1]{\par\needspace{4\baselineskip}\medskip\stepcounter{exer}\noindent{\color{navy}\bfseries\theexer.}\ \ {\small\estrelles{#1}}\ \ }

\newcommand{\seccio}[2]{\par\needspace{10\baselineskip}\bigskip
  \begin{tcolorbox}[colback=navy,colframe=navy,coltext=white,arc=2mm,boxrule=0pt,left=3mm,right=3mm,top=1.2mm,bottom=1.2mm]
  {\large\bfseries #1\quad #2}
  \end{tcolorbox}\vspace{-2pt}}

\newcolumntype{C}[1]{>{\centering\arraybackslash}p{#1}}

% quadre de recordatori (només a la fitxa d'exercicis)
\newcommand{\recorda}[1]{\par\smallskip\begin{tcolorbox}[colback=accent!7,colframe=accent!55,arc=2mm,boxrule=0.5pt,left=2.5mm,right=2.5mm,top=1mm,bottom=1mm]\small{\color{accent}\faLightbulb[regular]}\ \textbf{Recuerda:}\ #1\end{tcolorbox}}
% recta numèrica entera: \rectabuida{min}{max} i \recta{min}{max}{punts}
\newcommand{\rectabase}[2]{\draw[latex-latex,muted] (#1-0.8,0)--(#2+0.8,0);
  \foreach \i in {#1,...,#2}{\draw[muted] (\i,0.12)--(\i,-0.12) node[below=-1pt,font=\scriptsize,text=muted]{$\i$};}}
\newcommand{\rectabuida}[2]{\begin{tikzpicture}[x=0.78cm,baseline=0]\rectabase{#1}{#2}\end{tikzpicture}}
\newcommand{\recta}[3]{\begin{tikzpicture}[x=0.78cm,baseline=0]\rectabase{#1}{#2}
  \foreach \p in {#3}{\fill[accent] (\p,0) circle (0.11) node[above=1pt,font=\small\bfseries,text=navy]{$\p$};}\end{tikzpicture}}
\newcommand{\res}[1]{#1}

% =====================================================================
% Capçalera FentMates de les fitxes: un color per a cada curs.
% Els colors són els de la pàgina inicial (LEVELS a assets/js/materials.js);
% el logo ve de _eines/logo/pdf/ (python3 _eines/logo/genera.py).
%
% Ús, al final del preàmbul de la unitat:
%   % =====================================================================
% Capçalera FentMates de les fitxes: un color per a cada curs.
% Els colors són els de la pàgina inicial (LEVELS a assets/js/materials.js);
% el logo ve de _eines/logo/pdf/ (python3 _eines/logo/genera.py).
%
% Ús, al final del preàmbul de la unitat:
%   % =====================================================================
% Capçalera FentMates de les fitxes: un color per a cada curs.
% Els colors són els de la pàgina inicial (LEVELS a assets/js/materials.js);
% el logo ve de _eines/logo/pdf/ (python3 _eines/logo/genera.py).
%
% Ús, al final del preàmbul de la unitat:
%   \input{../comu/fentmates.tex}
%   \fmcurs{1eso}
% i, al començament del document:
%   \capcalera{Matemàtiques}{1r ESO}{Fitxa d'exercicis · Unitat 2}{Divisibilitat}
% Opcional abans de \input: \def\fmneed{14} (espai mínim abans de cada secció).
% =====================================================================
\makeatletter
\def\fm@logodir{../../logo/pdf}
\providecommand{\fmneed}{10}
\def\fm@cursnom{}

% clau del curs, color i color del text damunt del color
\newcommand{\fm@curs}[3]{\@namedef{fm@color@#1}{#2}\@namedef{fm@text@#1}{#3}}
\fm@curs{1eso}{E6B84A}{navy}
\fm@curs{2eso}{3B82F6}{white}
\fm@curs{3eso}{FF7A50}{white}
\fm@curs{4eso-a}{A78BFA}{white}
\fm@curs{4eso-b}{6D28D9}{white}
\fm@curs{1batx-mi}{9B1C1C}{white}
\fm@curs{1batx-ccss}{EF6B6B}{white}
\fm@curs{2batx-mii}{1E6B47}{white}
\fm@curs{2batx-ccss}{2BB3A3}{white}

\newcommand{\fmcurs}[1]{%
  \@ifundefined{fm@color@#1}{\PackageError{fentmates}{Curs desconegut: #1}{}}{}%
  \edef\fm@tmp{\noexpand\definecolor{curs}{HTML}{\@nameuse{fm@color@#1}}%
               \noexpand\colorlet{curstext}{\@nameuse{fm@text@#1}}}\fm@tmp
  \colorlet{cursfosc}{curs!72!black}%
  \def\fm@logo{\fm@logodir/vertical-#1.pdf}%
  \def\fm@marca{\fm@logodir/marca-#1.pdf}}

% ---------- capçalera de la primera pàgina ----------
% #1 matèria (petit), #2 curs, #3 tipus de document i unitat, #4 títol
\newcommand{\capcalera}[4]{%
  \gdef\fm@cursnom{#2}%
  \par\noindent
  \begin{tcolorbox}[colback=curs,colframe=curs,boxrule=0pt,arc=2.2mm,width=3.5cm,height=2.15cm,
      valign=center,left=3mm,right=2mm,top=1mm,bottom=1mm,nobeforeafter,box align=center]
    \color{curstext}{\footnotesize #1}\par\vspace{1pt}{\fontsize{19}{21}\selectfont\bfseries #2\par}%
  \end{tcolorbox}\hspace{5mm}%
  \begin{minipage}[c]{\dimexpr\linewidth-3.5cm-5mm-2.9cm\relax}
    {\color{cursfosc}\small\bfseries\MakeUppercase{#3}}\par\vspace{2pt}
    {\color{navy}\fontsize{20}{23}\selectfont\bfseries #4\par}
  \end{minipage}\hfill
  \begin{minipage}[c]{2.6cm}\raggedleft\includegraphics[width=2.6cm]{\fm@logo}\end{minipage}\par
  \vspace{2pt}}

% ---------- barres de secció del color del curs ----------
\renewcommand{\seccio}[2]{\par\needspace{\fmneed\baselineskip}\bigskip
  \begin{tcolorbox}[colback=curs,colframe=curs,coltext=curstext,arc=2mm,boxrule=0pt,left=3mm,right=3mm,top=1.2mm,bottom=1.2mm]
  {\large\bfseries #1\quad #2}
  \end{tcolorbox}\vspace{-2pt}}

% ---------- estrelles de dificultat del color del curs (a Batxillerat no se'n posen) ----------
\@ifundefined{estrelles}{}{%
  \renewcommand{\estrelles}[1]{{\color{curs}\ifcase#1\or\faStar\faStar[regular]\faStar[regular]\or\faStar\faStar\faStar[regular]\or\faStar\faStar\faStar\fi}}}

% ---------- capçalera i peu de les pàgines ----------
% primera pàgina: logos institucionals al peu
\fancypagestyle{primera}{\fancyhf{}\renewcommand{\headrulewidth}{0pt}
  \fancyfoot[L]{\raisebox{-0.32\height}{\includegraphics[height=0.78cm]{logo-gva.png}}\hspace{5mm}%
    \raisebox{-0.32\height}{\includegraphics[height=0.78cm]{logo-vermellar.png}}}
  \fancyfoot[R]{\scriptsize\color{muted}\href{https://creativecommons.org/licenses/by-nc-sa/4.0/deed.ca}{\ccbyncsa}\ \ CC BY-NC-SA 4.0 · fentmates.org\qquad\small\thepage}}
% la resta: el símbol a l'esquerra i el curs a la dreta, amb el títol al centre com abans
\fancyhead[L]{\raisebox{-0.3\height}{\includegraphics[height=0.62cm]{\fm@marca}}}
\fancyhead[R]{\small\bfseries\color{cursfosc}\fm@cursnom}
\renewcommand{\headrule}{\color{curs!45}\hrule height \headrulewidth}
\makeatother

%   \fmcurs{1eso}
% i, al començament del document:
%   \capcalera{Matemàtiques}{1r ESO}{Fitxa d'exercicis · Unitat 2}{Divisibilitat}
% Opcional abans de \input: \def\fmneed{14} (espai mínim abans de cada secció).
% =====================================================================
\makeatletter
\def\fm@logodir{../../logo/pdf}
\providecommand{\fmneed}{10}
\def\fm@cursnom{}

% clau del curs, color i color del text damunt del color
\newcommand{\fm@curs}[3]{\@namedef{fm@color@#1}{#2}\@namedef{fm@text@#1}{#3}}
\fm@curs{1eso}{E6B84A}{navy}
\fm@curs{2eso}{3B82F6}{white}
\fm@curs{3eso}{FF7A50}{white}
\fm@curs{4eso-a}{A78BFA}{white}
\fm@curs{4eso-b}{6D28D9}{white}
\fm@curs{1batx-mi}{9B1C1C}{white}
\fm@curs{1batx-ccss}{EF6B6B}{white}
\fm@curs{2batx-mii}{1E6B47}{white}
\fm@curs{2batx-ccss}{2BB3A3}{white}

\newcommand{\fmcurs}[1]{%
  \@ifundefined{fm@color@#1}{\PackageError{fentmates}{Curs desconegut: #1}{}}{}%
  \edef\fm@tmp{\noexpand\definecolor{curs}{HTML}{\@nameuse{fm@color@#1}}%
               \noexpand\colorlet{curstext}{\@nameuse{fm@text@#1}}}\fm@tmp
  \colorlet{cursfosc}{curs!72!black}%
  \def\fm@logo{\fm@logodir/vertical-#1.pdf}%
  \def\fm@marca{\fm@logodir/marca-#1.pdf}}

% ---------- capçalera de la primera pàgina ----------
% #1 matèria (petit), #2 curs, #3 tipus de document i unitat, #4 títol
\newcommand{\capcalera}[4]{%
  \gdef\fm@cursnom{#2}%
  \par\noindent
  \begin{tcolorbox}[colback=curs,colframe=curs,boxrule=0pt,arc=2.2mm,width=3.5cm,height=2.15cm,
      valign=center,left=3mm,right=2mm,top=1mm,bottom=1mm,nobeforeafter,box align=center]
    \color{curstext}{\footnotesize #1}\par\vspace{1pt}{\fontsize{19}{21}\selectfont\bfseries #2\par}%
  \end{tcolorbox}\hspace{5mm}%
  \begin{minipage}[c]{\dimexpr\linewidth-3.5cm-5mm-2.9cm\relax}
    {\color{cursfosc}\small\bfseries\MakeUppercase{#3}}\par\vspace{2pt}
    {\color{navy}\fontsize{20}{23}\selectfont\bfseries #4\par}
  \end{minipage}\hfill
  \begin{minipage}[c]{2.6cm}\raggedleft\includegraphics[width=2.6cm]{\fm@logo}\end{minipage}\par
  \vspace{2pt}}

% ---------- barres de secció del color del curs ----------
\renewcommand{\seccio}[2]{\par\needspace{\fmneed\baselineskip}\bigskip
  \begin{tcolorbox}[colback=curs,colframe=curs,coltext=curstext,arc=2mm,boxrule=0pt,left=3mm,right=3mm,top=1.2mm,bottom=1.2mm]
  {\large\bfseries #1\quad #2}
  \end{tcolorbox}\vspace{-2pt}}

% ---------- estrelles de dificultat del color del curs (a Batxillerat no se'n posen) ----------
\@ifundefined{estrelles}{}{%
  \renewcommand{\estrelles}[1]{{\color{curs}\ifcase#1\or\faStar\faStar[regular]\faStar[regular]\or\faStar\faStar\faStar[regular]\or\faStar\faStar\faStar\fi}}}

% ---------- capçalera i peu de les pàgines ----------
% primera pàgina: logos institucionals al peu
\fancypagestyle{primera}{\fancyhf{}\renewcommand{\headrulewidth}{0pt}
  \fancyfoot[L]{\raisebox{-0.32\height}{\includegraphics[height=0.78cm]{logo-gva.png}}\hspace{5mm}%
    \raisebox{-0.32\height}{\includegraphics[height=0.78cm]{logo-vermellar.png}}}
  \fancyfoot[R]{\scriptsize\color{muted}\href{https://creativecommons.org/licenses/by-nc-sa/4.0/deed.ca}{\ccbyncsa}\ \ CC BY-NC-SA 4.0 · fentmates.org\qquad\small\thepage}}
% la resta: el símbol a l'esquerra i el curs a la dreta, amb el títol al centre com abans
\fancyhead[L]{\raisebox{-0.3\height}{\includegraphics[height=0.62cm]{\fm@marca}}}
\fancyhead[R]{\small\bfseries\color{cursfosc}\fm@cursnom}
\renewcommand{\headrule}{\color{curs!45}\hrule height \headrulewidth}
\makeatother

%   \fmcurs{1eso}
% i, al començament del document:
%   \capcalera{Matemàtiques}{1r ESO}{Fitxa d'exercicis · Unitat 2}{Divisibilitat}
% Opcional abans de \input: \def\fmneed{14} (espai mínim abans de cada secció).
% =====================================================================
\makeatletter
\def\fm@logodir{../../logo/pdf}
\providecommand{\fmneed}{10}
\def\fm@cursnom{}

% clau del curs, color i color del text damunt del color
\newcommand{\fm@curs}[3]{\@namedef{fm@color@#1}{#2}\@namedef{fm@text@#1}{#3}}
\fm@curs{1eso}{E6B84A}{navy}
\fm@curs{2eso}{3B82F6}{white}
\fm@curs{3eso}{FF7A50}{white}
\fm@curs{4eso-a}{A78BFA}{white}
\fm@curs{4eso-b}{6D28D9}{white}
\fm@curs{1batx-mi}{9B1C1C}{white}
\fm@curs{1batx-ccss}{EF6B6B}{white}
\fm@curs{2batx-mii}{1E6B47}{white}
\fm@curs{2batx-ccss}{2BB3A3}{white}

\newcommand{\fmcurs}[1]{%
  \@ifundefined{fm@color@#1}{\PackageError{fentmates}{Curs desconegut: #1}{}}{}%
  \edef\fm@tmp{\noexpand\definecolor{curs}{HTML}{\@nameuse{fm@color@#1}}%
               \noexpand\colorlet{curstext}{\@nameuse{fm@text@#1}}}\fm@tmp
  \colorlet{cursfosc}{curs!72!black}%
  \def\fm@logo{\fm@logodir/vertical-#1.pdf}%
  \def\fm@marca{\fm@logodir/marca-#1.pdf}}

% ---------- capçalera de la primera pàgina ----------
% #1 matèria (petit), #2 curs, #3 tipus de document i unitat, #4 títol
\newcommand{\capcalera}[4]{%
  \gdef\fm@cursnom{#2}%
  \par\noindent
  \begin{tcolorbox}[colback=curs,colframe=curs,boxrule=0pt,arc=2.2mm,width=3.5cm,height=2.15cm,
      valign=center,left=3mm,right=2mm,top=1mm,bottom=1mm,nobeforeafter,box align=center]
    \color{curstext}{\footnotesize #1}\par\vspace{1pt}{\fontsize{19}{21}\selectfont\bfseries #2\par}%
  \end{tcolorbox}\hspace{5mm}%
  \begin{minipage}[c]{\dimexpr\linewidth-3.5cm-5mm-2.9cm\relax}
    {\color{cursfosc}\small\bfseries\MakeUppercase{#3}}\par\vspace{2pt}
    {\color{navy}\fontsize{20}{23}\selectfont\bfseries #4\par}
  \end{minipage}\hfill
  \begin{minipage}[c]{2.6cm}\raggedleft\includegraphics[width=2.6cm]{\fm@logo}\end{minipage}\par
  \vspace{2pt}}

% ---------- barres de secció del color del curs ----------
\renewcommand{\seccio}[2]{\par\needspace{\fmneed\baselineskip}\bigskip
  \begin{tcolorbox}[colback=curs,colframe=curs,coltext=curstext,arc=2mm,boxrule=0pt,left=3mm,right=3mm,top=1.2mm,bottom=1.2mm]
  {\large\bfseries #1\quad #2}
  \end{tcolorbox}\vspace{-2pt}}

% ---------- estrelles de dificultat del color del curs (a Batxillerat no se'n posen) ----------
\@ifundefined{estrelles}{}{%
  \renewcommand{\estrelles}[1]{{\color{curs}\ifcase#1\or\faStar\faStar[regular]\faStar[regular]\or\faStar\faStar\faStar[regular]\or\faStar\faStar\faStar\fi}}}

% ---------- capçalera i peu de les pàgines ----------
% primera pàgina: logos institucionals al peu
\fancypagestyle{primera}{\fancyhf{}\renewcommand{\headrulewidth}{0pt}
  \fancyfoot[L]{\raisebox{-0.32\height}{\includegraphics[height=0.78cm]{logo-gva.png}}\hspace{5mm}%
    \raisebox{-0.32\height}{\includegraphics[height=0.78cm]{logo-vermellar.png}}}
  \fancyfoot[R]{\scriptsize\color{muted}\href{https://creativecommons.org/licenses/by-nc-sa/4.0/deed.ca}{\ccbyncsa}\ \ CC BY-NC-SA 4.0 · fentmates.org\qquad\small\thepage}}
% la resta: el símbol a l'esquerra i el curs a la dreta, amb el títol al centre com abans
\fancyhead[L]{\raisebox{-0.3\height}{\includegraphics[height=0.62cm]{\fm@marca}}}
\fancyhead[R]{\small\bfseries\color{cursfosc}\fm@cursnom}
\renewcommand{\headrule}{\color{curs!45}\hrule height \headrulewidth}
\makeatother

\fmcurs{2eso}

\begin{document}
\thispagestyle{primera}

\capcalera{Matemáticas}{2\textsuperscript{o} ESO}{Ficha de ejercicios · Unidad 1}{Números naturales y enteros}
\vspace{10pt}
Nombre y apellidos: \rule{6.6cm}{0.4pt}\hfill Grupo: \rule{1.6cm}{0.4pt}\hfill Fecha: \rule{2.2cm}{0.4pt}
\vspace{8pt}

\begin{tcolorbox}[colback=soft,colframe=soft,arc=2mm,left=3mm,right=3mm,top=1.5mm,bottom=1.5mm]
\small Ejercicios graduados por niveles de dificultad:\quad \estrelles{1}\ básico \qquad \estrelles{2}\ medio \qquad \estrelles{3}\ avanzado.\\[2pt] Resuélvelos en el cuaderno mostrando todos los pasos. Haz los cálculos sin calculadora.
\end{tcolorbox}

% =====================================================
\seccio{1}{Los números naturales}

\ex{1} Calcula el cociente y el resto de cada división y haz la prueba ($D=d\cdot q+r$):\begin{tasks}(4)
  \task $247:12$
  \task $1000:37$
  \task $505:5$
  \task $89:90$
\end{tasks}

\ex{1} Indica si el resultado es un número natural. Si no lo es, explica por qué:\begin{tasks}(5)
  \task $7-12$
  \task $15\cdot 4$
  \task $20:6$
  \task $6:3+1$
  \task $3-3$
\end{tasks}

\ex{2} En una división el divisor es 15, el cociente 23 y el resto 9. ¿Cuál es el dividendo? ¿Podría ser 17 el resto de esta división?

\ex{2} Calcula mentalmente aplicando la propiedad distributiva:\begin{tasks}(4)
  \task $7\cdot 99$
  \task $12\cdot 15+12\cdot 5$
  \task $35\cdot 102$
  \task $8\cdot (10+5)$
\end{tasks}


% =====================================================
\seccio{2}{Múltiplos y divisores}

\ex{1} Escribe los cinco primeros múltiplos no nulos de:\begin{tasks}(3)
  \task $7$
  \task $12$
  \task $15$
\end{tasks}

\ex{1} Escribe todos los divisores de:\begin{tasks}(5)
  \task $24$
  \task $36$
  \task $45$
  \task $60$
  \task $17$
\end{tasks}

\ex{2} Di si son verdaderas o falsas y justifícalo:\begin{tasks}(1)
  \task 48 es múltiplo de 6.
  \task 7 es divisor de 77.
  \task 5 es múltiplo de 15.
  \task 1 es divisor de cualquier número.
  \task 0 es múltiplo de 9.
\end{tasks}


% =====================================================
\seccio{3}{Criterios de divisibilidad. Números primos}

\ex{2} Sin hacer la división, indica por cuáles de estos números es divisible cada uno: 2, 3, 5, 9, 10 y 11.\begin{tasks}(4)
  \task $1320$
  \task $4455$
  \task $7125$
  \task $2002$
\end{tasks}\recorda{Por el 11: resta la suma de las cifras de lugar impar y la de las de lugar par; tiene que dar 0 o un múltiplo de 11.}

\ex{2} ¿Qué cifras pueden ir en el lugar de $\square$?\begin{tasks}(1)
  \task $4\square7$ tiene que ser múltiplo de 3.
  \task $23\square$ tiene que ser múltiplo de 2 y de 5.
  \task $5\square2$ tiene que ser múltiplo de 9.
\end{tasks}

\ex{1} Clasifica en primos y compuestos. De los compuestos, escribe su descomposición:\begin{tasks}(8)
  \task $31$
  \task $51$
  \task $57$
  \task $61$
  \task $87$
  \task $91$
  \task $97$
  \task $1$
\end{tasks}


% =====================================================
\seccio{4}{Descomposición en factores primos}

\ex{1} Descompón en factores primos:\begin{tasks}(6)
  \task $60$
  \task $84$
  \task $126$
  \task $360$
  \task $1000$
  \task $1176$
\end{tasks}

\ex{2} ¿Qué número tiene cada descomposición?\begin{tasks}(3)
  \task $2^3\cdot3\cdot5^2$
  \task $2^2\cdot3^3$
  \task $3\cdot5\cdot7\cdot11$
\end{tasks}

\ex{3} ¿Por qué número, el más pequeño posible, hay que multiplicar 72 para que el resultado sea un cuadrado perfecto? ¿De qué número es el cuadrado?


% =====================================================
\seccio{5}{Mínimo común múltiplo y máximo común divisor}

\ex{1} Calcula listando múltiplos y divisores:\begin{tasks}(4)
  \task $\text{mcm}(12,18)$
  \task $\text{MCD}(12,18)$
  \task $\text{mcm}(8,20)$
  \task $\text{MCD}(8,20)$
\end{tasks}

\ex{2} Calcula el mcm y el MCD descomponiendo en factores primos:\begin{tasks}(3)
  \task $72$, $60$
  \task $84$, $126$
  \task $45$, $60$, $75$
  \task $140$, $300$
  \task $16$, $27$
\end{tasks}\recorda{\textbf{mcm}: factores comunes y no comunes con el mayor exponente. \textbf{MCD}: solo los comunes, con el menor exponente.}

\ex{2} Comprueba con 72 y 60 que el producto del MCD por el mcm es igual al producto de los dos números.

\ex{3} El MCD de dos números es 6 y el mcm, 90. Si uno de los números es 18, ¿cuál es el otro?

\ex{2} ¿Son primos entre sí estas parejas de números?\begin{tasks}(3)
  \task $15$ y $28$
  \task $21$ y $35$
  \task $9$ y $16$
\end{tasks}


% =====================================================
\seccio{6}{Problemas de mcm y MCD}

\ex{2} Dos líneas de autobús salen de la misma parada a las 7:00. La línea A pasa cada 15 minutos y la B, cada 20 minutos. ¿A qué hora volverán a coincidir?

\ex{2} Una floristería tiene 60 rosas y 84 claveles. Quiere hacer el mayor número posible de ramos iguales sin que sobre ninguna flor. ¿Cuántos ramos hará? ¿Cuántas rosas y cuántos claveles llevará cada ramo?

\ex{2} Queremos embaldosar el suelo de una habitación rectangular de 4,80 m × 3,60 m con baldosas cuadradas, las mayores posibles y sin cortar ninguna. ¿Cuánto debe medir el lado de cada baldosa? ¿Cuántas baldosas hacen falta?

\ex{2} Tres faros se encienden cada 8, 12 y 18 segundos, respectivamente. Si ahora se acaban de encender a la vez, ¿cuántos segundos pasarán hasta que vuelvan a encenderse juntos?

\ex{2} Un almacén tiene 150 libros de matemáticas y 210 de ciencias. Se quieren hacer paquetes iguales, con el máximo número de libros y con un solo tipo de libro en cada paquete. ¿Cuántos libros llevará cada paquete? ¿Cuántos paquetes se harán?

\ex{3} Si agrupamos los caramelos de una bolsa de 6 en 6, de 8 en 8 o de 10 en 10, siempre sobran 3. ¿Cuántos caramelos hay, como mínimo, en la bolsa (más de 3)?

\ex{3} Dos ruedas dentadas engranadas tienen 24 y 36 dientes. ¿Cuántas vueltas tiene que dar cada una para que vuelvan a estar en la posición inicial?


% =====================================================
\seccio{7}{Los números enteros}

\ex{1} Expresa con un número entero:\begin{tasks}(3)
  \task 3 grados bajo cero
  \task el segundo sótano
  \task una deuda de 50 €
  \task 1200 m de altitud
  \task el año 44 a. C.
  \task ganar 15 €
\end{tasks}

\ex{1} Escribe el valor absoluto y el opuesto de:\begin{tasks}(4)
  \task $-7$
  \task $12$
  \task $0$
  \task $-25$
\end{tasks}

\ex{1} Representa en la recta numérica y ordena de menor a mayor:\[-3,\quad 5,\quad {-8},\quad 0,\quad 2,\quad {-1},\quad {-12},\quad 7\]

\ex{2} Escribe todos los números enteros $x$ que cumplen:\begin{tasks}(4)
  \task $-4<x\le2$
  \task $\lvert x\rvert <3$
  \task $\lvert x\rvert =5$
  \task $-2\le x<-1$
\end{tasks}

\ex{2} Las temperaturas mínimas de una semana han sido: lunes $-3$\,°C, martes $-5$\,°C, miércoles $1$\,°C, jueves $0$\,°C, viernes $-2$\,°C, sábado $4$\,°C y domingo $-1$\,°C. ¿Qué día hizo más frío? Ordénalas de menor a mayor. ¿Qué diferencia hay entre la más alta y la más baja?


% =====================================================
\seccio{8}{Operaciones con enteros}

\ex{1} Calcula:\begin{tasks}(4)
  \task $(-7)+(-5)$
  \task $(-9)+4$
  \task $12+(-15)$
  \task $(-6)-(-10)$
  \task $8-(+13)$
  \task $(-3)-5$
  \task $0-(-7)$
  \task $(-11)+11$
\end{tasks}

\ex{1} Calcula agrupando los positivos y los negativos:\begin{tasks}(2)
  \task $5-8+3-7+2$
  \task $-4-6+10-3$
  \task $12-15-9+20$
  \task $-1+2-3+4-5+6$
\end{tasks}

\ex{2} Quita los paréntesis y los corchetes y calcula:\begin{tasks}(2)
  \task $7-(3-5+2)$
  \task $-(4-9)+(-2+6)$
  \task $10-\left[3-(5-8)\right]$
  \task $-\left[-(2-7)-(-3+1)\right]$
\end{tasks}\recorda{Un paréntesis precedido de un signo $-$ se quita cambiando el signo de todos los términos de dentro.}

\ex{1} Calcula aplicando la regla de los signos:\begin{tasks}(4)
  \task $(-6)\cdot (-4)$
  \task $7\cdot (-8)$
  \task $(-45):9$
  \task $(-56):(-7)$
  \task $0:(-5)$
  \task $(-2)\cdot (-3)\cdot (-5)$
  \task $(-1)\cdot (-1)\cdot (-1)\cdot (-1)$
  \task $72:(-8)$
\end{tasks}

\ex{2} ¿Qué número falta en cada caso?\begin{tasks}(3)
  \task $(-4)\cdot\square=28$
  \task $\square:(-6)=5$
  \task $\square\cdot(-9)=-81$
  \task $(-48):\square=-12$
  \task $(-5)+\square=-12$
  \task $\square-8=-3$
\end{tasks}

\ex{2} Un submarino está a 120 m de profundidad. Sube 45 m y después baja 80 m. ¿A qué profundidad se encuentra?

\ex{2} Pitágoras nació en el año 570 a. C. y murió en el año 495 a. C. ¿Cuántos años vivió?

\ex{3} Una cuenta bancaria tiene un saldo de 230 €. Se pagan tres recibos de 85 € cada uno y, después, se ingresan 120 €. ¿Cuál es el saldo final? ¿Ha estado la cuenta en negativo en algún momento?


% =====================================================
\seccio{9}{Operaciones combinadas}

\ex{1} Calcula respetando la jerarquía de las operaciones:\begin{tasks}(3)
  \task $5-3\cdot 4$
  \task $(5-3)\cdot 4$
  \task $18:(-3)+2\cdot (-5)$
  \task $-20:4\cdot (-2)$
  \task $3\cdot (-2)-4\cdot (-3)$
\end{tasks}\recorda{Primero, paréntesis; después, potencias y raíces; después, multiplicaciones y divisiones (de izquierda a derecha) y, al final, sumas y restas.}

\ex{2} Calcula:\begin{tasks}(2)
  \task $4-2\cdot \left[3-(6-10)\right]$
  \task $(-3)\cdot (2-7)+15:(-5)$
  \task $\left[(-8)+5\right]\cdot \left[6-(-2)\right]$
  \task $-6-\left[4-3\cdot (2-5)\right]$
  \task $25:(-5)-(-3)\cdot (4-6):2$
\end{tasks}

\ex{3} Un alumno ha hecho esta operación: \ $12-4:2\cdot(-3)=8:2\cdot(-3)=4\cdot(-3)=-12$. \ Explica qué error ha cometido y calcúlala bien.

\ex{3} Pon paréntesis para que las igualdades sean ciertas:\begin{tasks}(3)
  \task $6-4\cdot3+2=10$
  \task $6-4\cdot3+2=-14$
  \task $-8:2-6=2$
\end{tasks}


% =====================================================
\seccio{10}{Potencias y raíces con enteros}

\ex{1} Calcula:\begin{tasks}(4)
  \task $(-3)^{4}$
  \task $-3^{4}$
  \task $(-2)^{5}$
  \task $(-1)^{100}$
  \task $(-1)^{25}$
  \task $(-10)^{3}$
  \task $0^{7}$
  \task $(-7)^{0}$
\end{tasks}\recorda{Base negativa: exponente par $\to$ resultado positivo; exponente impar $\to$ resultado negativo. Cuidado: $-3^4=-(3^4)$, pero $(-3)^4=(-3)\cdot(-3)\cdot(-3)\cdot(-3)$.}

\ex{2} Sin hacer el cálculo, di qué signo tiene el resultado:\begin{tasks}(5)
  \task $(-7)^{12}$
  \task $(-3)^{15}$
  \task $-5^{4}$
  \task $-(-2)^{6}$
  \task $(-1)^{2026}$
\end{tasks}

\ex{2} Expresa como una sola potencia:\begin{tasks}(3)
  \task $(-2)^3\cdot(-2)^5$
  \task $(-5)^9:(-5)^6$
  \task $\left[(-3)^2\right]^4$
  \task $2^5\cdot(-3)^5$
  \task $(-12)^4:3^4$
  \task $7^3\cdot7^4:7^5$
\end{tasks}

\ex{2} Calcula aplicando las propiedades de las potencias:\begin{tasks}(4)
  \task $(-2)^7:(-2)^4$
  \task $\left[(-1)^3\right]^5\cdot(-1)^2$
  \task $5^4\cdot2^4$
  \task $(-6)^3:(-3)^3$
\end{tasks}

\ex{1} Calcula las raíces cuadradas (si existen):\begin{tasks}(6)
  \task $\sqrt{81}$
  \task $\sqrt{144}$
  \task $\sqrt{-25}$
  \task $\sqrt{0}$
  \task $\sqrt{400}$
  \task $\sqrt{-1}$
\end{tasks}

\ex{2} Calcula la raíz entera y el resto y haz la comprobación:\begin{tasks}(4)
  \task $\sqrt{50}$
  \task $\sqrt{90}$
  \task $\sqrt{130}$
  \task $\sqrt{200}$
\end{tasks}

\ex{2} Calcula:\begin{tasks}(2)
  \task $(-2)^{3}+3^{2}$
  \task $-2^{4}+(-2)^{4}$
  \task $5-2\cdot \sqrt{36}+(-1)^{7}$
  \task $(-3)^{2}\cdot 2-\sqrt{49}:(-7)$
  \task $\left[(-2)^{2}-3^{2}\right]^{2}:(-5)$
  \task $\sqrt{25-9}-(-3)^{3}:9$
\end{tasks}

\ex{3} Calcula paso a paso, respetando la jerarquía de las operaciones:\begin{tasks}(1)
  \task $\left[(-2)^{3}-(-3)^{2}\right]:(-17)+\sqrt{64}\cdot (-1)^{5}$
  \task $-3^{2}-\left[(-4)^{2}:(-2)^{3}-\sqrt{81}\right]\cdot (-2)$
  \task $\sqrt{5^{2}-3^{2}}\cdot (-2)^{3}-\left[(-6)^{2}:(-3)^{2}-2^{4}\right]^{2}:(-9)$
  \task $\left[2^{5}:((-2)^{2}-12)-(-1)^{20}\right]\cdot \sqrt{49}$
  \task $(-10)^{3}:(-5)^{2}+4\cdot \left[\sqrt{100}-(-3)^{3}:(-9)\right]^{2}$
\end{tasks}\recorda{Resuelve primero los paréntesis y corchetes de dentro hacia fuera. Dentro de cada uno, primero las potencias y las raíces, después las multiplicaciones y divisiones y, al final, las sumas y restas.}

\ex{3} Una parcela cuadrada tiene una superficie de 1225\,m². ¿Cuánto mide el lado? ¿Cuántos metros de valla hacen falta para rodearla?

\vfill
\begin{center}\color{muted}\small\faCheckCircle[regular]\ \ ¡Ánimo! Si te cuesta algún ejercicio, repasa los apuntes interactivos de la unidad.\end{center}

\end{document}
